\section{Knowledge Graph integration}
\label{sec:kg}

\subsection{Schema and backend}

The Phase 1/2 outputs already carry node and edge structure (\texttt{nodes\_station}, \texttt{nodes\_service}, \texttt{edges\_stops\_at}, \texttt{edges\_adjacent}, \texttt{nodes\_fault}). We promote them to a queryable graph database. The chosen backend is \textbf{Sparksee 5.2.3}~\cite{sparksee2024,martinez2017dex}, sourced from Maven Central (\texttt{com.sparsity:sparkseejava:5.2.3}); the jar bundles the Linux x86\_64 native libraries (\texttt{libsparksee.so}, \texttt{libsparkseejavawrap.so}, \texttt{libstlport.so}). We drive the JVM-side API from Python via JPype1 -- no system-level Java install is required because \texttt{scripts/install\_sparksee.sh} drops a portable Eclipse Temurin~17 JDK into \texttt{.tools/}. The schema is shown in \cref{lst:kg-schema}; the same DDL is portable to any labelled-property graph (Kuzu, Neo4j, Sparksee).

\begin{figure}[htbp]
\begin{lstlisting}[language=Cypher]
CREATE NODE TABLE Station(
    station_id            STRING,
    country               STRING,
    lat                   DOUBLE,
    lon                   DOUBLE,
    avg_historical_delay  DOUBLE,
    degree                INT64,
    PRIMARY KEY (station_id));
CREATE NODE TABLE TrainService(
    service_id        STRING, country STRING,
    train_class_code  INT64,  date DATE,
    is_disrupted      INT64,
    PRIMARY KEY (service_id));
CREATE NODE TABLE FaultEvent(
    fault_id     STRING, date DATE, description STRING,
    PRIMARY KEY (fault_id));
CREATE REL TABLE STOPS_AT(
    FROM TrainService TO Station,
    delay_minutes      DOUBLE, weather_severity   INT64,
    temperature DOUBLE, wind_speed DOUBLE,
    precipitation DOUBLE, snow_depth DOUBLE);
CREATE REL TABLE ADJACENT_TO(FROM Station TO Station, distance_km DOUBLE);
CREATE REL TABLE REPORTED_AT(FROM FaultEvent TO Station);
\end{lstlisting}
\caption{KG schema (Cypher DDL — the canonical, portable description). The Sparksee 5.2.3 backend creates the three node types and three edge types via the equivalent native API (\texttt{Graph.newNodeType}, \texttt{Graph.newEdgeType}, \texttt{Graph.newAttribute}) -- semantically identical, only the driver call sites differ.}
\label{lst:kg-schema}
\end{figure}

\subsection{Build statistics}

The Maven Central jar runs Sparksee in \emph{personal evaluation} mode unless the project's \texttt{sparksee.cfg} ships a licence string; eval mode caps each edge type at \SI{1}{M} edges. The full \texttt{STOPS\_AT} relation has \SI{16.59}{M} rows, so the build script auto-falls back to a representative \emph{hub subgraph} centred on the SHAP demo station \texttt{FI\_HKH}: every service that visits the hub or any of its 1-hop neighbours, plus the corresponding \texttt{STOPS\_AT}, \texttt{ADJACENT\_TO} and \texttt{REPORTED\_AT} edges. Holders of an academic / commercial licence can re-run with their licence string and load the full graph. \Cref{tab:kg-build} reports build wall-clock for the eval-mode hub subgraph on a Tesla~T4 / 4-core CPU box.

\begin{table}[htbp]
    \centering
    \caption{Sparksee 5.2.3 build statistics for the FI\_HKH hub subgraph on a Tesla~T4 / 4-core CPU box. JPype CRUD inserts at $\sim$\SI{18}{K} edges/s (single-threaded native API; the bulk \texttt{TypeLoader} would be faster but isn't reached here because the eval-mode cap binds first). Stage time is shared with the parquet feature store.}
    \label{tab:kg-build}
    \small
    \begin{tabular}{lrr}
        \toprule
        Step & Time & Rows \\
        \midrule
        Stage parquets (5 inputs concat + dedupe)        & $\sim$\SI{6}{min} & 16.59~M rows total \\
        Hub-subgraph trim (\texttt{station\_id} $\in$ FI\_HKH $\cup$ neighbours) & \SI{3}{s} & --- \\
        Schema creation (3 node + 3 edge types)          & \SI{0.1}{s} & --- \\
        \texttt{Station}            & \SI{0.0}{s}  & 3 \\
        \texttt{TrainService}       & \SI{1.0}{s}  & 16,434 \\
        \texttt{FaultEvent}         & \SI{0.0}{s}  & 0 \\
        \texttt{STOPS\_AT}          & \SI{5.6}{s} & 98,508 (\SI{18}{K}/s) \\
        \texttt{ADJACENT\_TO}       & \SI{0.0}{s} & 4 \\
        \texttt{REPORTED\_AT}       & \SI{0.0}{s} & 0 \\
        \midrule
        \textbf{Total Sparksee load} & $\sim$\textbf{7~s} & 98.5~K edges + 16.4~K nodes \\
        \bottomrule
    \end{tabular}
\end{table}

The on-disk \texttt{railway.gdb} file is~\SI{11}{MB} (hub subgraph) plus a \SI{640}{KB} OID index map.

\subsection{Bridge query: closing the loop from SHAP to KG}

The full point of the Knowledge Graph is to translate a \emph{black-box probability} into a \emph{queryable spatial-temporal neighbourhood}. After the SHAP report flags a high-risk stop, the analyst plugs $(\texttt{service\_id}, \texttt{station\_id}, \texttt{date})$ into the bridge query. \Cref{lst:kg-bridge} shows the canonical Cypher form (portable to Kuzu / Neo4j / Sparksee~6+ OpenCypher); \texttt{stop\_level/build\_kg.py::run\_bridge\_demo} implements an exactly-equivalent native-API translation against Sparksee 5.2.3 using \texttt{Graph.findEdge}, \texttt{Graph.explode} and \texttt{Graph.neighbors}.

\begin{figure}[htbp]
\begin{lstlisting}[language=Cypher]
MATCH (svc:TrainService {service_id: $sid})-[stop:STOPS_AT]->(st:Station {station_id: $stid})
OPTIONAL MATCH (f:FaultEvent)-[:REPORTED_AT]->(st)
  WHERE f.date >= date($d) - INTERVAL('14 days') AND f.date <= date($d)
OPTIONAL MATCH (st)-[adj:ADJACENT_TO]->(nbr:Station)
RETURN
  st.station_id          AS station,
  st.country             AS country,
  st.avg_historical_delay AS station_avg_delay,
  st.degree              AS degree,
  stop.delay_minutes     AS this_stop_delay,
  stop.weather_severity  AS weather_severity,
  count(DISTINCT f)      AS recent_faults_14d,
  count(DISTINCT nbr)    AS n_neighbours,
  collect(DISTINCT nbr.station_id) AS all_neighbours
\end{lstlisting}
\caption{KG bridge query: takes a SHAP-flagged $(\texttt{service\_id}, \texttt{station\_id}, \texttt{date})$ triple and returns the surrounding spatial-temporal context: recent faults, neighbouring stations, the actual stop delay and weather.}
\label{lst:kg-bridge}
\end{figure}

\paragraph{Live demo.} For the highest-confidence true-positive disruption from \texttt{xgb/B} (service \texttt{FI\_28\_2024-06-30} stopping at \texttt{FI\_HKH} on \texttt{2024-06-30}, $p=1.000$), the bridge query returns \cref{lst:kg-result} in \SI{7.7}{ms} on Sparksee 5.2.3 (vs the original Cypher form's \SI{52}{ms} on Kuzu) -- direct \texttt{findEdge} avoids the explode-and-filter pattern Cypher's planner falls back to.

\begin{figure}[htbp]
\begin{lstlisting}
{
  "station": "FI_HKH",       "country": "FI",
  "station_avg_delay": 6.84, "degree": 213,
  "this_stop_delay": 36.0,   "weather_severity": 0,
  "recent_faults_14d": 0,    "n_neighbours": 2,
  "sample_neighbours": ["FI_HVK", "FI_TKL"]
}
\end{lstlisting}
\caption{Bridge-query result for service \texttt{FI\_28\_2024-06-30} at station \texttt{FI\_HKH} on \texttt{2024-06-30}. The SHAP top driver \texttt{prev\_stop\_actual\_delay} is matched here by the actual \SI{36}{min} delay; \texttt{degree=213} confirms FI\_HKH is a hub.}
\label{lst:kg-result}
\end{figure}

\subsection{Reproducibility one-liner}

\begin{lstlisting}[language=bash]
# 1. Install Sparksee 5.2.3 (Maven Central jar) + portable Temurin JDK 17.
#    Idempotent — re-running is a no-op if both are already present.
source scripts/install_sparksee.sh

# 2. Build the KG and run the bridge demo. Drops any prior DB.
python stop_level/build_kg.py --rebuild
\end{lstlisting}

\paragraph{Licence tier.} The repo ships a \texttt{sparksee.cfg.example} with cache + storage tuning but no licence string. Holders of an academic / commercial Sparksee licence can copy it to \texttt{sparksee.cfg} and add their \texttt{sparksee.license = ...} line; the build then loads the full \SI{16.59}{M}-edge graph. Without a licence, the script auto-degrades to the FI\_HKH hub subgraph reported in \cref{tab:kg-build} -- still enough to demonstrate the bridge query end-to-end.
